% =====================================================================
%  CARNET D'ENTRAÎNEMENT — FICHE 5 (Bilans de matière)
%  THÈME 3 — Les proportions stœchiométriques (règle de trois)
%  Fiche TUTO
% =====================================================================
\documentclass[11pt,a4paper]{article}
\usepackage[utf8]{inputenc}
\usepackage[T1]{fontenc}
\usepackage{lmodern}
\usepackage[french]{babel}
\usepackage{amsmath,amssymb,mathtools}
\usepackage{icomma}
\usepackage[version=4]{mhchem}
\usepackage{siunitx}
\sisetup{output-decimal-marker={,},per-mode=symbol}
\DeclareSIUnit\Molar{\mole\per\litre}
\usepackage{xcolor}
\usepackage{colortbl}
\usepackage{tikz}
\usepackage[most]{tcolorbox}
\usepackage{enumitem}
\usepackage{array}
\usepackage{booktabs}
\usepackage{fancyhdr}
\usepackage[a4paper,top=1.6cm,bottom=1.5cm,left=1.9cm,right=1.9cm,headheight=24pt,headsep=6pt]{geometry}
\usetikzlibrary{arrows.meta,positioning,calc}

\definecolor{primary}{RGB}{142,93,168}
\definecolor{primarydark}{RGB}{82,45,108}
\definecolor{defcol}{RGB}{0,114,178}
\definecolor{thmcol}{RGB}{0,140,120}
\definecolor{methcol}{RGB}{90,90,100}
\definecolor{reacA}{RGB}{0,90,200}
\definecolor{reacB}{RGB}{210,30,40}
\definecolor{prodC}{RGB}{0,140,90}
\definecolor{prodD}{RGB}{198,93,9}
\newcommand{\Vm}{V_{\mathrm m}}
\newcommand{\nq}[1]{n_{\text{#1}}}

\pagestyle{fancy}\fancyhf{}
\fancyhead[L]{\small\textcolor{primarydark}{\textbf{Fiche 5} — Bilans de matière}}
\fancyhead[R]{\small\textcolor{primarydark}{\textsc{Tuto — Thème 3}}}
\fancyfoot[C]{\small\thepage}
\renewcommand{\headrulewidth}{0.8pt}
\renewcommand{\headrule}{\hbox to\headwidth{\color{primary}\leaders\hrule height \headrulewidth\hfill}}

\newtcolorbox{tutobox}[2][]{enhanced,breakable,boxrule=0.9pt,arc=2.5pt,
  left=3mm,right=3mm,top=2mm,bottom=2mm,colback=#2!6,colframe=#2,
  fonttitle=\bfseries,coltitle=white,title={#1}}

\setlength{\parindent}{0pt}
\setlength{\parskip}{2.5pt}

\begin{document}

\begin{center}
{\LARGE\bfseries\color{primarydark}Thème 3 — Les proportions stœchiométriques}\\[1pt]
{\color{primary}\rule{0.5\linewidth}{1.2pt}}\\[2pt]
{\itshape La « règle de trois » de la chimie : relier les espèces entre elles par les coefficients.}
\end{center}
\vspace{-3pt}

\begin{tutobox}[La double lecture des coefficients]{thmcol}
Pour $\ce{2 A + 3 B -> C + 6 D}$, les coefficients $(2,3,1,6)$ se lisent en \textbf{molécules}
\emph{ou} en \textbf{moles} : « 2~moles de \ce{A} réagissent avec 3~moles de \ce{B} pour donner
1~mole de \ce{C} et 6~moles de \ce{D} ». C'est la lecture \textbf{molaire} qui sert au chimiste.
\textbf{Ici, on suppose le rendement de \SI{100}{\percent}} (tout se transforme idéalement).
\end{tutobox}

\begin{tutobox}[La loi de proportionnalité]{defcol}
Les quantités \emph{consommées} et \emph{formées} sont \textbf{proportionnelles aux coefficients}.
Pour \ce{$a$A + $b$B -> $c$C + $d$D} :
\[
\boxed{\;\dfrac{|\Delta \nq{A}|}{a}=\dfrac{|\Delta \nq{B}|}{b}=\dfrac{\Delta \nq{C}}{c}=\dfrac{\Delta \nq{D}}{d}\;=\;\text{une même valeur (le « rapport de référence »)}}
\]
\textbf{Recette pratique.} Pour relier deux espèces \ce{X} et \ce{Y} :
$\;\boxed{\nq{Y}=\nq{X}\times\dfrac{\text{coeff.\ de Y}}{\text{coeff.\ de X}}}$ (règle de trois par les coefficients).
\end{tutobox}

\begin{tutobox}[La méthode « en escalier » (comme dans le livre)]{methcol}
On écrit l'équation, on porte les quantités connues sous chaque espèce, puis on relie par des flèches
$\times$(coefficient) pour \emph{monter} vers une espèce, $\div$(coefficient) pour \emph{descendre}.
\begin{center}
\begin{tikzpicture}[font=\footnotesize,scale=0.95]
  \def\cA{0}\def\cB{3.0}\def\cArr{5.6}\def\cC{7.7}\def\cD{10.2}
  \node[reacA,font=\bfseries] at (\cA,2.3) {2\,\ce{A}};
  \node at (\cA+1.3,2.3) {$+$};
  \node[reacB,font=\bfseries] at (\cB,2.3) {3\,\ce{B}};
  \draw[->,>=Stealth,line width=1pt] (\cArr-0.8,2.3) -- (\cArr+0.7,2.3);
  \node[prodC,font=\bfseries] at (\cC,2.3) {1\,\ce{C}};
  \node at (\cC+1.1,2.3) {$+$};
  \node[prodD,font=\bfseries] at (\cD,2.3) {6\,\ce{D}};
  \node[black!60] at (-1.7,1.5) {$t_0$};
  \node[reacA] at (\cA,1.5) {?};
  \node[reacB] at (\cB,1.5) {?};
  \node[prodC] at (\cC,1.5) {$0$};
  \node[prodD] at (\cD,1.5) {$0$};
  \node[black!60] at (-1.7,0.4) {$t_f$};
  \node[reacA] at (\cA,0.4) {$0$};
  \node[reacB] at (\cB,0.4) {$0$};
  \node[prodC,font=\bfseries] at (\cC,0.4) {$0{,}100$};
  \node[prodD] at (\cD,0.4) {?};
  \draw[->,>=Stealth,primary,line width=0.9pt] (\cC,0.7) to[bend left=14] node[above,black,font=\scriptsize]{$\times 2$} (\cA,1.2);
  \draw[->,>=Stealth,primary,line width=0.9pt] (\cC,0.7) to[bend left=10] node[above,pos=0.6,black,font=\scriptsize]{$\times 3$} (\cB,1.2);
  \draw[->,>=Stealth,primary,line width=0.9pt] (\cC,0.35) to[bend right=18] node[below,black,font=\scriptsize]{$\times 6$} (\cD,0.35);
  \draw[primary!30,rounded corners=3pt] (-2.3,-0.1) rectangle (11.4,2.7);
\end{tikzpicture}
\end{center}
On part de la quantité \emph{visée} de \ce{C} ($0{,}100$~mol) et on remonte :
$\nq{A}=2\times0{,}100=0{,}200$, $\nq{B}=3\times0{,}100=0{,}300$, $\nq{D}=6\times0{,}100=0{,}600$~mol.
\end{tutobox}

\begin{tutobox}[La chaîne complète d'un calcul de proportion]{primary}
\textbf{1.} Pondérer l'équation (Thème 1). \quad
\textbf{2.} Convertir la donnée en \textbf{moles} (Thème 2). \quad
\textbf{3.} Appliquer la \textbf{règle de trois} par les coefficients pour l'espèce cherchée. \quad
\textbf{4.} \textbf{Reconvertir} le résultat dans l'unité demandée ($m=nM$, $V=n\Vm$, $C=n/V$).\\[2pt]
\emph{Exemple express.} Combustion \ce{CH4 + 2 O2 -> CO2 + 2 H2O}, on veut \SI{88}{\gram} de \ce{CO2}
($M=44$) : $\nq{CO2}=88/44=2$~mol $\Rightarrow \nq{CH4}=2\times\frac{1}{1}=2$~mol et
$\nq{O2}=2\times\frac{2}{1}=4$~mol.
\end{tutobox}

\end{document}
